% Generated draft: requires mathematical and source review.
\documentclass[11pt]{article}
\usepackage[margin=25mm]{geometry}
\usepackage{fontspec}
\setmainfont{Latin Modern Roman}
\usepackage{amsmath,amssymb,mathtools,unicode-math}
\setmathfont{Latin Modern Math}
\usepackage{enumitem,needspace,xurl,hyperref}
\setlength{\parindent}{0pt}
\setlength{\parskip}{5pt}
\setlength{\emergencystretch}{3em}
\newcommand{\N}{\mathbb{N}}
\newcommand{\Z}{\mathbb{Z}}
\newcommand{\Q}{\mathbb{Q}}
\newcommand{\R}{\mathbb{R}}
\newcommand{\C}{\mathbb{C}}
\newcommand{\F}{\mathbb{F}}
\newcommand{\A}{\mathbb{A}}
\newcommand{\PP}{\mathbb{P}}
\newcommand{\E}{\mathbb{E}}
\newcommand{\eps}{\varepsilon}
\newcommand{\dd}{\,\mathrm{d}}
\begin{document}
\section*{UnsolvedMath ID: 11000029}
\textbf{Source problem code:} catalogue entry
\subsection{Definitions and mathematical statement}
Let \(G=\operatorname{Mod}_g\) in the source's surface convention, and fix a word metric from a finite generating set. Given a specified family \(\mathscr G\) of geodesics, a natural relative quasiconvexity condition for \(H\leq G\) is the existence of \(R\geq0\) such that every selected geodesic with endpoints in \(H\) lies in
\[
 N_R(H)=\{x:d(x,H)\leq R\}.
\]
The family must at least supply a geodesic for each relevant pair of endpoints; otherwise the condition could be vacuous.

The problem asks to characterize the subgroups of \(\operatorname{Mod}_g\) that are quasiconvex with respect to \emph{some collection of geodesics}. The excerpt does not specify whether the family is a fixed global combing, may depend on \(H\), has equivariance or consistency axioms, or requires all geodesics rather than one chosen geodesic per pair. These choices and the metric must be fixed before the predicate is unique; the displayed condition provides explicit notation, not a silent selection among them.

The source distinguishes this from convex cocompactness as defined in [FMo], whose usual geometric formulation involves quasiconvexity of an orbit in Teichm\"uller space. It states that an orbit embedding \(\operatorname{Mod}_g\to\operatorname{Teich}_g\) is exponentially distorted, citing [FLMi], so the two ambient geodesic problems cannot simply be identified. Precise metric, genus and theorem hypotheses are not supplied.
\subsection{Short English statement}
Characterize mapping-class subgroups that stay near a suitably chosen family of geodesics, keeping this word-geometric question distinct from Teichm\"uller convex cocompactness.
\subsection{Sources}
\begin{itemize}
\item \emph{Problems on Mapping Class Groups and Related Topics}, edited by Farb, Problem 3.8, supplied PDF page 34. \url{https://www.math.uchicago.edu/~farb/papers/mcgbook.pdf}. Recorded access 2026-07-29; extraction pdf-font-structure.
\item Source references [FMo] for convex cocompactness and [FLMi] for exponential distortion. Full bibliography entries and theorem hypotheses are absent.
\item Local UnsolvedMath registry, \texttt{data/unsolvedmath/problems.json}, permanent ID \(11000029\), exact source code \texttt{catalogue entry}; supplied OPEN-TRIAGE summary reports related convex-cocompact partial work without an exact resolving citation.
\end{itemize}
\end{document}
